\section{Empirical Strategy}\label{sec:strategy}

The strategy exploits province-level meat-price variation as an
external salience treatment.
The treatment is measured by the National Bureau of Statistics
and merged from outside the CFPS survey, so the household's
current expectation never appears on both sides of any
regression.
The primary test is a reduced-form regression of forecast errors
on the meat shock.
Two decomposition equations identify the channels through which
the forecast error arises.

\subsection{Treatment Variable}\label{sec:meatshock}

The treatment variable is the cumulative log change in provincial
meat CPI between two consecutive CFPS waves, defined in
equation~\eqref{eq:meat_shock_data}.
Three properties matter for identification.
First, it is external to the survey: it contains no information
derived from respondents.
Second, it is cumulative: a household surveyed in 2020
experienced the full path of meat-price movements since the
2018 wave.
Third, it varies at the province level, which allows me to
absorb region-by-wave fixed effects and still retain
within-region identifying variation.

The cumulative measure is the natural choice given the CFPS
survey design.
CFPS is fielded every two years, with interviews spanning
several months within each wave.
Households report expectations once per wave, and those reports
reflect whatever price information they accumulated since the
prior interview.
The cumulative log change between wave midpoints therefore
captures the total price signal that a household plausibly
processed.

Two measurement considerations deserve attention.
First, the NBS meat CPI index combines pork, beef, mutton, and
poultry, with pork receiving the dominant weight.
During the ASF episode, pork drove the bulk of the meat CPI
movement; the index therefore tracks the pork shock with high
fidelity during the identification window.
Second, the cumulative measure aggregates monthly observations
equally.
If households weight recent months more heavily, the treatment
variable understates effective salience, attenuating the
estimated coefficient toward zero.

The 2018--2019 ASF episode provides the sharpest source of
variation.
Provincial herd losses varied because culling policy, transport
restrictions, and initial herd density differed across provinces.
This heterogeneity generates cross-sectional dispersion in the
treatment that is plausibly orthogonal to province-level demand
conditions.

\subsection{Three-Equation System}\label{sec:three_equations}

Primary test: forecast errors (Equation~3).
The reduced-form implication of overreaction is that households
in high-shock provinces subsequently overpredict inflation:
\begin{equation}\label{eq:eq3}
  \mathrm{FE}_{ipt}
  = \alpha
  + \beta_3\,\mathrm{MeatShock}_{p,\,t-1\to t}
  + \gamma' X_{ipt}
  + \delta_{r(p)\times t}
  + \varepsilon_{ipt},
\end{equation}
where $\mathrm{FE}_{ipt}
  \equiv \pi^{\mathrm{headline}}_{p,\,t\to t+1}
  - \tilde{\mu}_{ipt}$ is realised provincial headline inflation
over the next inter-wave period minus the household's directional
expectation scaled to CPI units.
A negative forecast error means that the household expected more
inflation than materialised.
Under overreaction, $\beta_3 < 0$: high-shock provinces generate
more negative forecast errors.

The forecast-error equation occupies the primary position because
it requires the fewest auxiliary assumptions.
The sign of $\beta_3$ alone tests overreaction.
If households responded rationally, their expectations would
adjust by no more than the shock's actual headline CPI weight,
and forecast errors would be bounded by
$\omega \approx 0.05$ regardless of the household's belief about
shock persistence (Section~\ref{sec:framework}).
Any $|\hat{\beta}_3|$ exceeding this bound rejects rationality
without requiring knowledge of the conversion factor between
ordinal expectations and CPI units.

The identification of $\beta_3$ comes from within-region
variation in meat-price shocks.
Two households in different provinces within the same region and
survey wave face identical aggregate conditions but different
meat-price histories.
The coefficient $\beta_3$ measures whether the household in
the higher-shock province produces a more negative forecast
error.

Decomposition: expectations (Equation~1).
\begin{equation}\label{eq:eq1}
  \mu_{ipt}
  = \alpha
  + \beta_1\,\mathrm{MeatShock}_{p,\,t-1\to t}
  + \gamma' X_{ipt}
  + \delta_{r(p)\times t}
  + \varepsilon_{ipt},
\end{equation}
where $\mu_{ipt} \in \{-1,0,1\}$ is the ordinal inflation
expectation.
A positive $\hat{\beta}_1$ indicates that higher meat-price
shocks shift expectations upward.

Equation~1 identifies whether the shock moves beliefs.
If households are inattentive or if meat prices lack
salience, $\beta_1$ would be zero.
Because the dependent variable is ordinal, the coefficient
measures shifts in the probability of reporting ``up'' versus
``same'' or ``down,'' not a cardinal change in expected
inflation.
The ordinal scale limits what $\beta_1$ can tell us about the
magnitude of the expectation response in CPI units, which is
precisely why the forecast-error equation is the primary test.

Decomposition: pass-through (Equation~2).
\begin{equation}\label{eq:eq2}
  \pi^{\mathrm{headline}}_{p,\,t\to t+1}
  = \alpha
  + \beta_2\,\mathrm{MeatShock}_{p,\,t-1\to t}
  + \delta_{r(p)\times t}
  + u_{pt},
\end{equation}
where $\pi^{\mathrm{headline}}_{p,\,t\to t+1}$ is realised
province-level headline CPI inflation over the subsequent
inter-wave period.
This regression runs at the province-wave level (93 observations).
If $\beta_2 \leq 0$, the shock does not raise future headline
inflation, closing the second half of the overreaction logic.

The pass-through equation pins down what a rational household
should have expected.
A zero or negative $\beta_2$ means the shock reverses. Provinces
that received large meat-price increases subsequently saw lower,
not higher, headline inflation.
The reversal follows ASF supply-shock dynamics.
As herds recovered through 2021--2022, pork supply expanded and
meat prices fell, dragging down headline CPI in the provinces
that experienced the sharpest earlier increases.
The 93-observation sample size motivates the wild cluster
bootstrap described in Section~\ref{sec:inference}.

\subsubsection{Decomposition Logic and the Overreaction Test}%
\label{sec:overreaction_decomposition}

The three coefficients are linked by an approximate
decomposition:
$\beta_3 \approx \beta_2 - k\,\beta_1$, where $k$ converts
the ordinal expectation into implied inflation units.
This identity clarifies the overreaction test.

The forecast error is realised inflation minus expected inflation.
Regressing realised inflation on the meat shock gives $\beta_2$.
Regressing expected inflation on the meat shock gives
$k\,\beta_1$.
If expectations respond ($\beta_1 > 0$) but realised CPI does
not ($\beta_2 \leq 0$), the forecast error is necessarily
negative because the household raised its expectation in response
to a signal that did not raise the object it was forecasting.

The sign test $\beta_3 < 0$ does not require knowledge of $k$.
The CFPS expectation question is ordinal, so the mapping from
$\mu \in \{-1,0,1\}$ to CPI units is unknown.
But the forecast error is constructed in CPI units, so
$\beta_3$ is directly interpretable.
Whatever $k$ is, if the forecast error is systematically more
negative in high-shock provinces, households in those provinces
overpredicted inflation.

The decomposition also clarifies why $|\hat{\beta}_3|$
exceeds the rational bound by such a wide margin.
Under rational expectations, $\beta_1$ should equal
$\omega\,\hat{\rho} / k$, where $\omega \approx 0.05$ is meat's
CPI weight and $\hat{\rho}$ is the household's persistence
belief.
Even if $\hat{\rho} = 1$, the rational forecast-error coefficient
is bounded by $|\beta_3^R| \leq 0.05$.
Finding $|\hat{\beta}_3| \approx 6$ means the actual expectation
response far exceeds the rational response.

\subsection{Identification Assumptions}%
\label{sec:identification_assumptions}

The strategy requires two conditions.

The first condition is relevance: the meat-price shock must
predict household forecast errors.
Pork is widely consumed, purchased frequently, and priced
visibly; ASF generated extensive media coverage that was
difficult to miss even for inattentive consumers.
Three features support relevance: pork accounts for roughly
60 percent of China's meat consumption and is purchased several
times per week; the ASF episode received extensive state and
social media coverage; and prices at wet markets are posted on
visible boards and updated daily.

The second condition is conditional exogeneity: the meat-price
shock must be uncorrelated with province-level
determinants of forecast errors other than through salience,
conditional on fixed effects and controls.
Region-by-wave fixed effects absorb regional business cycles,
monetary policy transmission, and nationwide shocks hitting
all provinces within a region simultaneously.
The remaining identifying variation is within-region
cross-province differences in meat-price shocks over time.

The main threat is a province-specific demand shock that
simultaneously raises meat prices and expectations through
a real-income channel.
Two features mitigate this concern.
First, ASF provides a biologically determined supply shock.
The disease entered China through contaminated feed imports
and spread along pig transportation corridors, with provincial
infection severity determined by initial herd density and
proximity to early outbreak sites.
Second, the commodity-placebo test in
Section~\ref{sec:mechanism} shows that grain shocks of
comparable magnitude do not produce the same forecast-error
pattern.
If local demand were driving both prices and expectations,
grain prices should respond to the same demand shock and
produce a similar expectation effect.

A second concern is that province-level meat-price variation
might proxy for local food-price inflation more broadly.
Region-by-wave fixed effects partially address this by
absorbing region-level food-price trends.
The within-region variation in meat prices is largely
idiosyncratic to the ASF episode and does not coincide with
similar dispersion in non-meat food categories.

The design also requires stable unit treatment: household $i$'s
forecast error depends
on the meat shock in province $p$ where the household resides,
not on meat shocks in other provinces.
Region-by-wave fixed effects absorb any national-level or
region-level price signal, so the identifying variation reflects
local price conditions.
The CFPS expectation question asks about general price changes,
not pork prices specifically; households that follow national
pork-price news would report higher expectations in all
provinces, which would be absorbed by the wave fixed effect.

What the fixed effects absorb and what they leave.
Region-by-wave fixed effects partition the 31 provinces into
six NBS macro-regions and absorb any time-varying shock common
to provinces within a region.
This removes aggregate monetary policy shocks, regional business
cycles, and region-specific food supply disruptions.
The remaining variation is within-region cross-province
differences in meat-price shocks.
The fixed effects do not absorb province-specific shocks that
are correlated with the meat shock.
If a province experienced an unobserved demand boom that
raised both meat prices and expectations, this would bias
$\beta_3$ toward zero rather than away from zero.
The direction of potential omitted-variable bias therefore
works against the overreaction finding.

\subsection{Inference}\label{sec:inference}

With 31 province clusters, asymptotic cluster-robust standard
errors can over-reject \citep{Cameron2008}.
Simulation evidence suggests that 30--50 clusters is a
borderline case where asymptotic approximations may or may not
perform well.

I report wild cluster bootstrap $p$-values throughout, using
the \citet{Webb2023} six-point weight distribution with 9,999
replications.
The six-point distribution is preferred over the Rademacher
two-point distribution because it provides better size control
for cluster counts between 10 and 50, precisely the range that
applies here.
The bootstrap is implemented under the null hypothesis of
zero treatment effect (restricted residuals).

Point estimates and conventional cluster-robust standard errors
are reported alongside bootstrap $p$-values.
All significance statements in the text refer to bootstrap
$p$-values unless otherwise noted.
